\documentclass[english, parskip=full]{scrartcl}
\usepackage[utf8]{inputenc}  % Kodierung der Datei
\usepackage[english]{babel}  % Sprache auf Deutsch 
\usepackage{color}
\usepackage{graphicx, gensymb}
\usepackage{amsfonts, cancel, amssymb}
\usepackage{amsmath, amsthm, array, esint}
\usepackage{booktabs, multirow}  % schönere Tabellen
\usepackage{caption}  % Anpassung der Captions
\usepackage{scrlayer-scrpage}    % Manipulation des Seitenstils
\pagestyle{scrheadings}  % Stil für die Seite setzen
\automark{section}  % Abschnittsnamen als Seitenbeschriftung 
\setheadsepline{.5pt}    % Kopzeile durch Linie abtrennen
\usepackage[hidelinks]{hyperref}  % Links und weitere PDF-Features
\usepackage{typearea, tikz, pgffor, verbatim}
\usetikzlibrary{calc, arrows.meta,arrows, graphs, quotes, positioning}
\usepackage[cache=false, outputdir=build]{minted}
\usemintedstyle{vs}
\usepackage{placeins} % provides \FloatBarrier

\definecolor{bg}{rgb}{0.9,0.9,0.9}
\newcommand\numberthis{\addtocounter{equation}{1}\tag{\theequation}}
\usepackage{svg}
\usepackage{siunitx}
\usepackage{physics}
\usepackage{tensor}
\renewcommand{\bra}{}
\newcommand{\kom}{}
\newcommand{\brak}{}
\renewcommand{\braket}{}
\newcommand{\intx}{}
\newcommand{\intr}{}
\newcommand{\kla}{}
\newcommand{\klam}{}
\newcommand{\klammer}{}
\newcommand{\oper}{}
\newcommand{\tecxt}{}
\newcommand{\Bra}{}
\newcommand{\Ket}{}
\renewcommand{\i}{\text{i}}
\newcommand{\e}{\text{e}}
\newcommand*{\Fourier}{\textsc{Fourier}\ }
\newcommand*{\F}{\mathcal{F}}
\newcommand*{\sinc}{\text{sinc\,}}
\newcommand{\conv}{\mathop{\scalebox{3.5}{\raisebox{-0.38ex}{\makebox[0.5\width][c]{$\ast$}}}}\limits}

\newcommand*{\titel}{02 Time-dependent Ginzburg-Landau Equation}
\newcommand*{\autor}{Joachim Schwardt}

\hypersetup{pdfauthor={\autor}, pdftitle={\titel}}

%\titlehead{titlehead}
\subject{Theory of Superconductivity}
\title{\titel}
\author{\autor}
\date{\today}

\begin{document}

%\begin{titlepage}
\maketitle  % Titel setzen
%\tableofcontents  % Inhaltsverzeichnis setzen
%\end{titlepage}

Consider the Lagrange density
\begin{align}
\mathcal{L} &= \i\hbar \psi^\ast \partial_t\psi - \alpha |\psi|^2 - \frac{\beta}{2} |\psi|^4 - \frac{\hbar^2}{2m^\ast} |\nabla\psi|^2.
\end{align}
The action is defined as $\mathcal{S} := \int \mathcal{L}(\psi, \partial_\mu\psi) \,\mathrm{d}^4x$.
Assuming the variational principle $\delta\mathcal{S} = 0$ one finds the Euler-Lagrange equations for fields.
We will derive these using $\psi \rightarrow \psi + \delta\psi$
\begin{align}
\mathcal{S} &\rightarrow \mathcal{S}' = \intx\int_{ }^{ } \mathcal{L}\kla\left( \psi + \delta\psi, \partial_\mu \psi + \partial_\mu \delta_\psi \right) \, \mathrm{d}^4x = \\
&= \intx\int_{ }^{ } \mathcal{L} + \frac{\partial \mathcal{L}}{\partial \psi} \delta\psi + \frac{\partial \mathcal{L}}{\partial \partial_\mu \psi} \partial_\mu \delta\psi \, \mathrm{d}^4x + \mathcal{O}(\delta\psi)^2 = \\
&= \mathcal{S} + \intx\int_{ }^{ } \klam\left[ \frac{\partial \mathcal{L}}{\partial \psi} - \partial_\mu \frac{\partial \mathcal{L}}{\partial \partial_\mu \psi}  \right] \delta\psi \, \mathrm{d}^4x + \klam\left[ \frac{\partial \mathcal{L}}{\partial \partial_\mu \psi} \delta\psi \right]_{t_i}^{t_f}.
\end{align}
We assume the variation $\delta\psi$ to vanish for the initial and final time, which removes the term leftover from partial integration.
Then using $\mathcal{S}' - \mathcal{S} = 0$ we find
\begin{align}
\frac{\partial \mathcal{L}}{\partial \psi} &= \partial_\mu \frac{\partial \mathcal{L}}{\partial \partial_\mu \psi}.
\end{align}
In our particular case, for the variation with respect to $\psi^\ast$, this yields
\begin{align}
0 &= \partial_\mu \frac{\partial \mathcal{L}}{\partial \partial_\mu \psi^\ast} - \frac{\partial \mathcal{L}}{\partial \psi^\ast} = \\
&= \nabla\cdot \kla\left( -\frac{\hbar^2}{2m^\ast} \nabla\psi \right) - \i\hbar\partial_t \psi + \alpha\psi + \beta|\psi|^2\psi,\\
\Rightarrow\quad \i\hbar\partial_t\psi &= -\frac{\hbar^2}{2m^\ast} \nabla^2\psi + \alpha\psi + \beta|\psi|^2 \psi.
\end{align}
In the uniform case we may assume $\nabla\psi = 0$ and split of the equilibrium value using $\psi(t) = \sqrt{\frac{-\alpha}{\beta}} f(t)$.
Then we find
\begin{align}
\i\hbar f'(t) &= \alpha f(t) - \alpha |f(t)|^2 f(t).
\end{align}
Introducing the time scale $\tau := \frac{\hbar}{-\alpha} > 0$ and expressing $f$ using modulus and phase via $f(t) = g(t) \e^{\i\phi(t)}$ gives
\begin{align}
\i\tau g' \e^{\i\phi} + \i\tau \i\phi' g\e^{\i\phi} &= -g\e^{\i\phi} + g^3\e^{\i\phi},\\
\Rightarrow\quad \i\tau g' - \phi' g &= -g + g^3.
\end{align}
Since $g,\phi$ are real functions, the imaginary and real parts of this equation must vanish individually.
The former gives $\i\tau g'(t) = 0$ and therefore tells us that $g = \tecxt\text{const}.$
The latter leads to
\begin{align}
\phi'(t) g &= g \kla\left( 1 - g^2 \right).
\end{align}
Consider suddenly switching on superconducting pairing at $t=0$, i.e. $\psi(t=0) = 0$.
This is equivalent to $g \equiv 0$ since $g$ is conserved, and thus $\psi \equiv 0$ is the only solution.
\\
Now consider a jump in the parameter $\alpha$, described by the initial condition $\psi(t=0) = \sqrt{\frac{-\tilde{\alpha}}{\beta}}$ with $\tilde{\alpha} < 0$.
Then we find $g \equiv \sqrt{\frac{-\tilde{\alpha}}{\beta}}$ and may solve the equation for $\phi'$
\begin{align}
\phi(t) &= \phi(0) + \kla\left( 1 - g^2 \right) t = \kla\left( 1 + \frac{\tilde{\alpha}}{\beta} \right)t.
\end{align}
This gives
\begin{align}
\psi(t) &= \sqrt{\frac{-\tilde{\alpha}}{\beta}} \e^{\i \kla\left( 1 + \frac{\tilde{\alpha}}{\beta} \right) t},
\end{align}
which describes an oscillation of the complex phase with frequency $\nu := 1 + \frac{\tilde{\alpha}}{\beta}$.



%\begin{thebibliography}{9}
%\bibitem{example} description, \url{https://www.exmapleurl}, day.\,month~2021.
%\end{thebibliography}

\end{document}